\section{Simulation Results}
\label{sec:results}
All 170 measured executions complete the frozen workload and pass the
independent raw-evidence audit. This section separates protection traffic,
recovery preparation and whole-workflow completion. Figures show all five
paired admission orders and 95\% intervals where applicable.
\begin{table}[t]
\centering
\caption{Base-setting absolute endpoints (seconds), averaged over five seeds.
These descriptive means do not replace paired penalty contrasts.}
\label{tab:endpoints}
\setlength{\tabcolsep}{3pt}
\begin{tabular}{@{}lrrr@{}}
\toprule
Arm & F0 WCT & F1 WCT & F1 batch \\
\midrule
Unprotected & 41.565 & 42.686 & 441.871 \\
DRAM-only & 41.565 & 42.453 & 439.506 \\
Full AegisKV & 41.488 & 42.407 & 439.038 \\
\bottomrule
\end{tabular}
\end{table}


\subsection{Healthy protection cost}
Figure~\ref{fig:healthy} shows substantial copying with little resolved
WCT change. Full transfers 5,497 GiB per workload, versus
2,867 GiB for DRAM-only (91.8\% more). Its paired mean
WCT difference is -0.078 s (95\% CI [-0.229, +0.073] s);
unprotected has the same healthy control times. The interval includes zero.
Background copies can overlap compute, while memory and closed-loop batching
change schedules. These modeled timings exclude real controller/kernel costs.

\begin{figure*}[t]
\centering\includegraphics[width=\textwidth]{figures/aligned-healthy.pdf}
\caption{Healthy execution. (a) Full minus control mean WCT, aggregated within
each seed before comparison. (b) Bytes actually transferred for protection
across the complete workload, including work on later-aborted copies.
Traffic is in TiB ($2^{40}$ bytes). Open circle, square, up/down triangle
and plus markers identify seeds 0--4; filled diamonds and bars show means
and 95\% Student-$t$ intervals. Modeled zero latency cost would not imply zero physical
controller or kernel overhead.}
\Description{Paired healthy workflow-latency differences and protection
traffic for unprotected, DRAM-only and full treatments.}
\label{fig:healthy}
\end{figure*}

\subsection{Failure impact on complete workflows}
The raw ECDFs largely overlap (Figure~\ref{fig:failure} and
Table~\ref{tab:endpoints}). Relative to DRAM-only, full's mean failure-penalty
contrast is +0.032 s (95\% CI [-0.159, +0.222] s);
the focus contrast is +1.661 s (95\% CI [-2.931, +6.254] s).
Against unprotected recomputation, the whole-workload contrast is
-0.201 s (95\% CI [-0.669, +0.267] s).
None establishes a stable advantage across these admission orders.

\begin{figure*}[t]
\centering\includegraphics[width=\textwidth]{figures/aligned-failure.pdf}
\caption{Base-setting failure outcomes. (a) Descriptive WCT ECDFs over 160
workflows in each of five runs per arm; the 800 values are not independent
replicates. (b,c) Paired F1-minus-F0 penalty contrasts for the whole-workload
mean and the prespecified focus workflow. Negative values favor full.
Intervals use five admission orders, not workflow-level pseudo-replicates.}
\Description{Raw workflow completion distributions and controlled contrasts
against unprotected recomputation and DRAM-only protection.}
\label{fig:failure}
\end{figure*}

\subsection{Recovery state and preparation}
Full copies HBM for four of five focus requests; one uses DRAM alone
(Figure~\ref{fig:recovery}). In seed 0, both arms restore 36,864 prompt tokens,
but full reaches prefix-ready 0.711 s sooner. The corresponding workflow
penalty contrast is still +0.029 s.
In the largest adverse focus contrast (seed 4), full's workflow penalty
exceeds DRAM-only by 8.278 s despite faster
focus-request recovery. The trace shows a different branch sharing the
recovery worker becoming the last branch before the join; the final request
also waits longer. This post-hoc example explains why a restored prefix or
faster preparation cannot establish workflow-level benefit.

\begin{figure*}[t]
\centering\includegraphics[width=\textwidth]{figures/aligned-recovery.pdf}
\caption{The five prespecified focus requests. (a) DRAM/full bars partition each
prompt into HBM copy, DRAM load and recomputation. (b) Paired prefix-ready
times include queueing and preparation; lines connect treatments within a
seed. (c) Prefix-ready saving against the F0-corrected focus-WCT penalty
contrast; numbers identify seeds (1--3 nearly coincide). Seed 4 prepares
faster yet finishes its workflow later. No regression is fitted.}
\Description{Per-seed recovery-source proportions, paired prefix preparation times,
and their relationship to corrected workflow completion penalties.}
\label{fig:recovery}
\end{figure*}

\subsection{Mechanism ablations}
Every whole-workload and focus interval in Figure~\ref{fig:ablations}
includes zero; the zero-Slack focus contrast is exactly zero in all five
seeds. There is no resolved latency benefit for an individual mechanism here.
Retention nevertheless changes healthy traffic: removing Slack raises it to
6,873 GiB and LRU to 7,175 GiB, versus
5,497 GiB for full. Fixing the initial quota uses
5,107 GiB. These byte costs are mechanism observations,
not evidence of a failure-latency or hardware-throughput improvement.

\begin{figure*}[t]
\centering\includegraphics[width=\textwidth]{figures/aligned-ablations.pdf}
\caption{Ablation minus full: (a,b) failure penalties corrected by each arm's
own F0 run; (c) healthy copy traffic. Positive values favor full. Slack
sets the signal to zero; AIMD fixes the initial quota; LRU replaces retention.
All keep request ownership, copying and cleanup. Each estimate uses five
paired admission orders; the byte effect does not establish a WCT gain.}
\Description{Whole-workload, focus-workflow and healthy-byte effects of removing
Slack, fixing quota control, and replacing retention with LRU.}
\label{fig:ablations}
\end{figure*}

\subsection{Sensitivity to modeled conditions}
All fourteen whole-workload/focus intervals across the seven settings
include zero (Figure~\ref{fig:sensitivity}). At 5 GB/s, the focus estimate is
-2.239 s (95\% CI [-12.889, +8.410] s),
showing substantial uncertainty. With two sequential faults, the whole-workload
contrast is +0.046 s (95\% CI [-0.249, +0.341] s).
Service composition and admission order can change signs and dispersion;
these experiments support a conditional mechanism study, not a universal
performance ranking.

\begin{figure*}[t]
\centering\includegraphics[width=\textwidth]{figures/aligned-sensitivity.pdf}
\caption{Full minus DRAM failure-penalty contrasts across the declared
settings. $C$ is workflow concurrency and $\beta$ the serialized fraction of
modeled batch service. Two faults use their two prespecified focus workflows
and reuse base F0 controls. Open markers show all seeds; intervals are exploratory,
unadjusted comparisons. Negative values favor full.}
\Description{Paired sensitivity to concurrency, batch service composition,
path bandwidth and two sequential failures.}
\label{fig:sensitivity}
\end{figure*}
